\section{Sharpness, repair taxonomy, and strict Route verdict}
\label{sec:sharpness-route}

The all-prime-square quantifier is not decorative.  Every finite truncation
admits an exact periodic witness, and the empty truncation requires a separate
argument.  These controls delimit the theorem and prevent a finite
computation from being promoted to an all-prime statement.

\subsection{Every finite prime-square source has periodic points}

Let \(P_0\subset\Pset\) be finite and define
\begin{equation}
  Q=\prod_{p\in P_0}p^2.
  \label{eq:finite-Q}
\end{equation}

\begin{proposition}[Finite-exclusion sharpness]
\label{prop:finite-sharpness}
The subshift obtained by enforcing admissibility only for the primes in
\(P_0\) fails the full source's periodic-collapse and proximality
conclusions.
\end{proposition}

\begin{proof}
Assume first that \(P_0\neq\varnothing\).  Define
\begin{equation}
  x_n=1
  \quad\Longleftrightarrow\quad
  n\equiv1\pmod Q.
  \label{eq:finite-witness}
\end{equation}
For every \(p\in P_0\), divisibility \(p^2\mid Q\) gives
\[
  \supp(x)\bmod p^2=\{1\bmod p^2\},
\]
so \(x\) satisfies every retained admissibility constraint.  It is nonzero
and \(Q\)-periodic.  Moreover, its least positive period is \(Q\): if
\(d>0\) is a period, then \(x_1=1\) forces \(x_{1+d}=1\), hence
\(Q\mid d\).  Since a nonempty \(P_0\) gives \(Q\geq4\), this is a
nontrivial periodic orbit.

If \(P_0=\varnothing\), then \(Q=1\) and the constraints are vacuous.
The points \(0^{\mathbb Z}\) and \(1^{\mathbb Z}\) are distinct fixed
points, so the system is not proximal and does not have a singleton periodic
ledger.  This separate empty-set case is essential: one cannot call the
period-one all-ones point a nontrivial \(Q\)-cycle.
\end{proof}

For the concrete singleton \(P_0=\{2\}\), the point
\((0111)^{\mathbb Z}\) supplies an additional least-period-four witness.  Its
support modulo four is \(\{1,2,3\}\), so it omits residue zero.  This example
is illustrative only; \cref{prop:finite-sharpness} is proved for every finite
\(P_0\), including the empty set.

\begin{figure}[t]
\centering
\begin{tikzpicture}[
  node distance=8mm and 10mm,
  box/.style={draw,rounded corners=2pt,align=center,text width=0.255\textwidth,
    minimum height=25mm,inner sep=4pt,font=\small},
  arrow/.style={-{Latex[length=2.4mm]},thick}
]
\node[box,draw=proofgreen,fill=proofgreen!6] (infinite)
  {\textbf{All primes}\\fresh \(p^2\) for every window\\proximality and
  singleton periodic ledger};
\node[box,draw=stopred,fill=stopred!5,right=of infinite] (finite)
  {\textbf{Finite nonempty \(P_0\)}\\
  \(Q=\prod_{p\in P_0}p^2\)\\
  \(x_n=1\iff n\equiv1\pmod Q\)\\least period \(Q\)};
\node[box,draw=stopred,fill=stopred!5,right=of finite] (empty)
  {\textbf{Empty \(P_0\)}\\constraints vacuous\\
  \(0^{\mathbb Z}\neq1^{\mathbb Z}\)\\two fixed points};

\draw[arrow,stopred] (infinite) -- (finite);
\draw[arrow,stopred] (finite) -- (empty);
\node[font=\scriptsize,align=center,fill=white,inner sep=1.5pt]
  at ($(infinite)!0.5!(finite)+(0,18mm)$)
  {delete infinitely\\many exclusions};
\node[font=\scriptsize,align=center,fill=white,inner sep=1.5pt]
  at ($(finite)!0.5!(empty)+(0,18mm)$)
  {delete the rest};

\node[draw=mutedpurple,fill=mutedpurple!6,rounded corners=2pt,
  below=10mm of finite,text width=0.74\textwidth,align=center,inner sep=5pt,
  font=\small] (repairs)
  {Products/extensions import cycles; inducing changes time; nonfactor maps
  delete hypotheses; an aperiodic zeta changes the observable.  None is a
  same-contract repair.};
\draw[arrow,mutedpurple,dashed] (finite) -- (repairs);
\draw[arrow,mutedpurple,dashed] (empty) -- (repairs);
\end{tikzpicture}
\caption{Sharpness of the all-prime source.  The empty finite set is a
separate edge case, not a nontrivial \(Q\)-cycle.}
\label{fig:sharpness}
\end{figure}


\subsection{Apparent repairs and their type changes}

\begin{table}[t]
\centering
\small
\caption{Why common repairs do not remain in the frozen contract.}
\label{tab:repairs}
\begin{tabularx}{\textwidth}{@{}>{\raggedright\arraybackslash}X
  >{\raggedright\arraybackslash}Xl@{}}
\toprule
Proposed repair & Mathematical effect & Verdict \\
\midrule
Retain finitely many prime-square exclusions
  & admits the exact periodic points in \cref{prop:finite-sharpness}
  & source change \\
Take a product with a periodic system
  & imports an external periodic ledger through an extension
  & direction error \\
Induce on or return to a subset
  & changes time, marker, and primitive orbit type
  & clock change \\
Use a noncontinuous, nononto, or nonequivariant map
  & leaves the factor category used by the proof
  & scope deletion \\
Replace fixed-point counts by an aperiodic or measure zeta
  & changes the determinant owner
  & observable change \\
Call the powers \(z^r\) new rational-prime primitives
  & converts traversals into atoms
  & repetition/type error \\
Call \([1]\) a full source transfer operator
  & assigns a ledger package to an unsupported state space
  & ownership error \\
\bottomrule
\end{tabularx}
\end{table}

The one-point factor and the identity factor are lawful positive controls.
The former realizes the theorem in its most collapsed form; the latter
retains the exact source.  Neither creates a nontrivial primitive ledger.

\subsection{Strict Route-A record}

The exact factor theorem does not promote the candidate through the broader
arithmetic program.  Its strict Route-A tuple is
\[
\begin{gathered}
  (\mathrm{A0\_FAIL},\ \mathrm{A1\_FAIL},\
   \mathrm{A2\_ANALYTIC\_DETERMINANT},\\
   \mathrm{A3\_FAIL},\ \mathrm{A4\_FAIL}).
\end{gathered}
\]
The coordinates have distinct meanings:

\begin{itemize}
\item \textbf{A0 fails.} Every rational-prime-square exclusion is explicitly inserted
  into the source grammar.  No rational-prime primitives emerge endogenously.
\item \textbf{A1 fails.} The primitive ledger is proved exactly, including temporal
  repetition, but contains only one trivial fixed orbit and has no
  rational-prime primitive support.
\item \textbf{A2 passes only at the analytic-determinant label.} The exact inverse
  determinant \(1-z\) is a polynomial and is realized by \([1]\), but it
  records only the singleton ledger.
\item \textbf{A3 fails.} There is no completed arithmetic divisor, gamma factor,
  functional equation, \(T\log T\) growth, explicit formula, or natural
  same-ledger Weil compression.
\item \textbf{A4 fails.} No full-state transfer family, fixed self-adjoint
  Hilbert--Polya operator, same-clock prime trace identity, or completed target
  divisor is defined.
\end{itemize}

Thus the overall verdict is \texttt{ROUTE\_A\_REJECTED}.  Route B invocation
is false because the same-object primitive ledger and completed analytic
structure already fail.  The terminal statements are
\texttt{STOP\_PROXIMAL\_PERIODIC\_RIGIDITY},
\texttt{STOP\_SINGLETON\_PRIMITIVE\_SUPPORT}, and
\texttt{STOP\_TRIVIAL\_ONE\_MINUS\_Z\_DIVISOR}.  These conclusions are
retrospective renderings of the proved theorem and frozen contract, not
prospective predictions.
